\chapter{总结}
溯界实现了以任务契约、来源用途和效果级授权为核心的智能体工具防护原型，提供13类业务配置及任务、审批、来源和影响查询。系统支持真实MCP调用、跨角色委派、跨会话记忆、结构化返回快照、版本化规则初筛、敏感写入硬拒绝与账本溯源展示。

提交程序固定于已验证基线。Ubuntu、Windows和macOS分别有1107、1106和1107项回归通过；四业务确定性运行在提交包解压后完成16个单元，离线规则评测可按随附样例重算。历史48单元比较另列源版本并保留对照组危险效果。以上证据支持当前原型的工程行为复核，尚不能给出真实模型攻击成功率或通用防护结论。

后续验证需在同一模型与代码版本下比较业务完成、危险效果和代价，补充未见任务、有效消融与重复实验，并完成完整宿主/工具拓扑的隔离、跨机部署及性能测量。作品以源码哈希、实际状态和独立核验结果组织证据，便于评审复现已完成部分并识别尚待验证的边界。
